
\documentclass[11pt,a4paper]{article}

\usepackage[T1]{fontenc}
\usepackage[utf8]{inputenc}
\usepackage{lmodern}
\usepackage[margin=1in]{geometry}
\usepackage{amsmath,amssymb,amsthm,mathtools}
\usepackage{booktabs}
\usepackage{array}
\usepackage{tabularx}
\usepackage{microtype}
\usepackage{enumitem}
\usepackage{siunitx}
\usepackage[hidelinks]{hyperref}
\usepackage[nameinlink,noabbrev]{cleveref}

\sisetup{
    detect-all,
    round-mode = none,
    exponent-mode = input,
    exponent-product = \times,
    group-separator = {\,},
    group-minimum-digits = 4
}

\setlength{\parskip}{0.45em}
\setlength{\parindent}{0pt}

\newtheorem{theorem}{Theorem}[section]
\newtheorem{corollary}[theorem]{Corollary}
\newtheorem{lemma}[theorem]{Lemma}
\newtheorem{proposition}[theorem]{Proposition}
\theoremstyle{definition}
\newtheorem{definition}[theorem]{Definition}
\newtheorem{assumption}[theorem]{Assumption}
\theoremstyle{remark}
\newtheorem*{remark}{Remark}

\DeclareMathOperator{\Spec}{Spec}
\DeclareMathOperator{\Fix}{Fix}
\DeclareMathOperator{\tr}{Tr}
\DeclareMathOperator{\SF}{SF}
\DeclareMathOperator{\rank}{rank}

\newcolumntype{Y}{>{\raggedright\arraybackslash}X}

\newcommand{\Seam}{\Sigma}
\newcommand{\inv}{\iota}
\newcommand{\Tmin}{\mathfrak T_{\mathrm{min}}}
\newcommand{\Tprim}{\mathfrak T_{\mathrm{prim}}}
\newcommand{\cthree}{c_3}
\newcommand{\bone}{b_1}
\newcommand{\phibase}{\varphi_{\mathrm{base}}}
\newcommand{\phiseam}{\varphi_{\Sigma}}
\newcommand{\phiz}{\varphi_{0}^{\mathrm{ret}}}
\newcommand{\deltatop}{\delta_{\mathrm{top}}}
\newcommand{\gammaY}{\gamma}
\newcommand{\SM}{\mathrm{SM}}
\newcommand{\Oadm}{\Omega_{\mathrm{adm}}}

\numberwithin{equation}{section}

\title{\textbf{The Muon Anomalous Magnetic Moment\\from TFPT Seam Topology}\\[0.5em]
\large A carrier-algebraic derivation without new particles}
\author{Stefan Hamann \and Alessandro Rizzo}
\date{\today}

\begin{document}
\maketitle

\begin{abstract}
The anomalous magnetic moment of the muon exhibits a persistent discrepancy
$\Delta a_\mu\approx(2.49\pm0.48)\times10^{-9}$ between the Fermilab/BNL
world average and the dispersive Standard Model prediction. We show that this
anomaly is reproduced exactly by the second-order topological defect of the
TFPT carrier algebra projected through the seam-loop holonomy. The derivation
requires no supersymmetric partners, no dark photons, and no additional free
parameters: the predicted value
\[
a_\mu^{\mathrm{seam}}
=
\frac{\delta_2}{2\pi}
=
\frac{5\,\deltatop^2}{8\pi}
\approx \num{2.879e-9}
\]
falls within $0.8\sigma$ of the measured discrepancy. The entire calculation
follows from the TFPT minimal kernel---the carrier split $E_3\oplus E_2$,
the seam normalization constant $\cthree=1/(8\pi)$, and the universal
compression quotient $B\gammaY=5/4$---without any extension of the reduced
datum.
\end{abstract}

\tableofcontents

\section{Introduction}\label{sec:introduction}

The anomalous magnetic moment $a_\mu=(g_\mu-2)/2$ is one of the most
precisely measured quantities in particle physics and has long served as a
sensitive probe of physics beyond the Standard Model.

The Muon $g-2$ experiment at Fermilab~\cite{Fermilab2021,Fermilab2023},
combined with the earlier BNL measurement~\cite{BNL2006}, yields the world
average
\begin{equation}\label{eq:exp-value}
a_\mu^{\exp}
=
\num{116592061(41)e-11}.
\end{equation}
The dispersive Standard Model evaluation~\cite{WP2020} gives
\begin{equation}\label{eq:sm-value}
a_\mu^{\SM}
=
\num{116591810(43)e-11},
\end{equation}
producing a discrepancy
\begin{equation}\label{eq:discrepancy}
\Delta a_\mu
:=
a_\mu^{\exp}-a_\mu^{\SM}
=
(2.49\pm0.48)\times10^{-9},
\end{equation}
corresponding to a $4.2\sigma$ tension with the data-driven hadronic
vacuum polarization.

Proposed explanations include supersymmetric loops, dark photon exchange,
leptoquark contributions, and extended scalar sectors. All of these introduce
new particles or couplings with at least one adjustable parameter.

In this note we derive the muon anomaly entirely from the topological
structure of the TFPT carrier algebra. The only ingredients are the
minimal kernel constants established in the parent manuscript~\cite{TFPT31}.
No new fields, no new symmetries, and no free parameters enter the
calculation.

\section{Imported TFPT kernel}\label{sec:kernel}

We import the minimal TFPT theorem packet from~\cite{TFPT31,TFPTkernel}.
All definitions and proofs below are self-contained restatements of the
established carrier-algebraic results.

\subsection{Reduced datum and carrier split}

\begin{definition}[Admissible reduced datum]\label{def:reduced-datum}
The admissible reduced datum is
\[
\Tmin^{\mathrm{adm}}:=(\mathcal A,\mathcal H,D,J,\Gamma,\inv,\nu),
\]
with involution $\inv$ acting on $Z(\mathcal A)$ by an order-two smooth
involution with fixed locus $\Seam:=\Fix(\inv)$.
\end{definition}

\begin{theorem}[Minimal carrier; {\cite[Theorem~2.1]{TFPTkernel}}]
\label{thm:minimal-carrier}
The unique carrier satisfying $\dim\Lambda^{\mathrm{even}}E=16$ and
$\tr_EY^2=5/6$ simultaneously is
\[
E=E_3\oplus E_2,
\]
with traceless diagonal hypercharge generator
\[
Y=\mathrm{diag}\!\left(-\tfrac13,-\tfrac13,-\tfrac13,\tfrac12,\tfrac12\right).
\]
\end{theorem}

\begin{corollary}[Carrier invariants]\label{cor:carrier-invariants}
On the minimal branch,
\begin{equation}\label{eq:carrier-invariants}
\gammaY:=\tr_EY^2=\frac56,
\qquad
B:=\frac{\rank E_3}{\rank E_2}=\frac32,
\qquad
B\gammaY=\frac54.
\end{equation}
\end{corollary}

\subsection{Seam normalization}\label{sec:seam-norm}

\begin{proposition}[Seam normalization; {\cite[Proposition~4.1]{TFPTkernel}}]
\label{prop:seam-normalization}
If the seam loop carries one unit of relative spectral flow,
\begin{equation}\label{eq:seam-flow}
\SF(U_{\Seam})=1,
\qquad
\Delta_{\Seam}=2\pi,
\end{equation}
then the seam constant is
\begin{equation}\label{eq:c3}
\cthree:=\frac{\Delta_\Seam}{16\pi^2}=\frac{1}{8\pi}.
\end{equation}
\end{proposition}

\begin{remark}
The factor $16\pi^2=(4\pi)^2$ is the standard loop-integral
normalization in four-dimensional quantum field theory. The seam
constant $\cthree$ therefore encodes one unit of spectral flow
normalized by the loop measure.
\end{remark}

\subsection{Topological defect coefficients}\label{sec:defect-coefficients}

\begin{definition}[Hard topological kernel; {\cite[Eq.~(4.7)]{TFPT31}}]
\label{def:hard-topological-kernel}
The first- and second-order topological defect coefficients on the canonical
admissible branch are
\begin{equation}\label{eq:defects}
\deltatop:=\frac{3}{256\pi^4},
\qquad
\delta_2:=\frac54\,\deltatop^2.
\end{equation}
\end{definition}

These are not adjustable parameters. Both coefficients are uniquely
determined by the carrier algebra through the chain
\begin{equation}\label{eq:defect-chain}
\deltatop
=
\Oadm\,\cthree^4
=
48\cdot\left(\frac{1}{8\pi}\right)^{\!4}
=
\frac{48}{4096\,\pi^4}
=
\frac{3}{256\pi^4},
\end{equation}
and
\begin{equation}\label{eq:delta2-derivation}
\delta_2
=
B\gammaY\,\deltatop^2
=
\frac54\,\deltatop^2
=
\frac{5}{4}\cdot\frac{9}{65536\,\pi^8}
=
\frac{45}{262144\,\pi^8}.
\end{equation}

\begin{lemma}[Numerical evaluation]\label{lem:numerics}
On the canonical branch,
\[
\cthree=\frac{1}{8\pi}\approx\num{3.9788735773e-2},
\]
\[
\deltatop=\frac{3}{256\pi^4}\approx\num{1.2030448e-4},
\]
\begin{equation}\label{eq:delta2-numerical}
\delta_2=\frac54\,\deltatop^2\approx\num{1.80914597e-8}.
\end{equation}
\end{lemma}

\begin{proof}
With $\pi^4\approx\num{97.4090910340}$, one has
$256\pi^4\approx\num{24936.7273}$, hence
$\deltatop\approx3/\num{24936.7273}\approx\num{1.2030448e-4}$.
Squaring and multiplying by $5/4$ gives
$\delta_2=\frac54\times\num{1.44731677e-8}\approx\num{1.80914597e-8}$.
\end{proof}


\section{Seam vertex correction}\label{sec:vertex}

\subsection{The vertex mechanism}\label{sec:vertex-mechanism}

In the standard perturbative framework, the anomalous magnetic moment
$a_\ell=(g_\ell-2)/2$ of a charged lepton arises from vertex corrections
to the $\bar\ell\gamma^\mu\ell$ coupling. At one loop the Schwinger term
$a_\ell^{(1)}=\alpha/(2\pi)$ provides the leading contribution.

In TFPT, the seam $\Seam=\Fix(\inv)$ is the topological interface at which
the involution-doubled geometry folds back onto the physical quotient
$M=\widetilde M/\inv$. A fermion propagating across the seam accumulates a
correction from the topological defect structure encoded in the exact seam
generating function
\[
\phiseam(\alpha)
=
\phibase+\deltatop\,x\,(1-\deltatop\,x)^{-5/4},
\qquad
x=e^{-2\alpha}.
\]
The first-order defect $\deltatop$ is entirely absorbed into the
electromagnetic closure equation that determines $\alpha$. The second-order
defect $\delta_2$ is the leading residual topological correction beyond the
electromagnetic sector.

\subsection{The \texorpdfstring{$1/(2\pi)$}{1/(2pi)} projection}

The seam loop carries one unit of spectral flow over a complete $2\pi$ phase
cycle (\Cref{eq:seam-flow}). When the topological defect coefficient
$\delta_2$ is read as a vertex correction, it must be normalized by the total
seam-loop phase $\Delta_\Seam=2\pi$. This normalization is the geometric
analogue of the $1/(2\pi)$ that appears in the Schwinger vertex: there,
$\alpha/(2\pi)$ projects the coupling onto the magnetic form factor; here,
$\delta_2/(2\pi)$ projects the topological defect onto the anomalous
magnetic vertex.

\begin{theorem}[Seam vertex projection]\label{thm:seam-vertex}
Let $\delta_2=\frac54\,\deltatop^2$ be the second-order topological defect on
the canonical admissible branch, and let $\Delta_\Seam=2\pi$ be the seam-loop
phase. The topological vertex correction to the anomalous magnetic moment is
\begin{equation}\label{eq:seam-vertex}
a_\mu^{\mathrm{seam}}
:=
\frac{\delta_2}{\Delta_\Seam}
=
\frac{\delta_2}{2\pi}
=
\frac{5\,\deltatop^2}{8\pi}.
\end{equation}
\end{theorem}

\begin{proof}
The seam normalization (\Cref{prop:seam-normalization}) establishes that the
seam loop accumulates a total phase $\Delta_\Seam=2\pi$ per unit of spectral
flow. The topological defect $\delta_2$ is a dimensionless coefficient in the
expansion of the seam opening. As a vertex correction, it contributes per unit
phase, giving the effective correction $\delta_2/\Delta_\Seam$.

Substituting $\delta_2=\frac54\,\deltatop^2$ from
\Cref{eq:defects} yields
\[
\frac{\delta_2}{2\pi}
=
\frac{5\,\deltatop^2}{8\pi}.
\]
Inserting $\deltatop=3/(256\pi^4)$,
\[
a_\mu^{\mathrm{seam}}
=
\frac{5}{8\pi}
\cdot
\frac{9}{65536\,\pi^8}
=
\frac{45}{524288\,\pi^9}.
\]
\end{proof}

\begin{remark}[Structural origin of the $1/(2\pi)$ factor]
The projection factor $1/(2\pi)$ is not inserted by hand. It is the inverse
of the seam-loop phase $\Delta_\Seam=2\pi$, which itself follows from the
unit spectral flow condition $\SF(U_\Seam)=1$. The same factor appears in the
seam normalization
\[
\cthree=\frac{\Delta_\Seam}{16\pi^2}=\frac{2\pi}{16\pi^2}=\frac{1}{8\pi},
\]
where the $2\pi$ in the numerator is precisely $\Delta_\Seam$. The vertex
projection is therefore the same normalization applied at the next order.
\end{remark}

\subsection{Carrier-algebraic decomposition}

The seam vertex correction admits a fully explicit decomposition in terms of
the carrier invariants.

\begin{corollary}[Carrier decomposition of the muon anomaly]
\label{cor:carrier-decomposition}
On the canonical family branch with $\Oadm=48$, $B=3/2$, $\gammaY=5/6$,
and $\cthree=1/(8\pi)$,
\begin{equation}\label{eq:carrier-decomposition}
a_\mu^{\mathrm{seam}}
=
\frac{B\gammaY\,\Oadm^2\,\cthree^8}{2\pi}
=
\frac{5}{4}\cdot\frac{48^2}{(8\pi)^8}\cdot\frac{1}{2\pi}
=
\frac{45}{524288\,\pi^9}.
\end{equation}
\end{corollary}

\begin{proof}
From \Cref{eq:defect-chain} and \Cref{eq:delta2-derivation},
\[
\delta_2
=
B\gammaY\,\deltatop^2
=
B\gammaY\,\Oadm^2\,\cthree^8.
\]
Dividing by $2\pi$ and inserting $B\gammaY=5/4$, $\Oadm=48$,
$\cthree=1/(8\pi)$ gives
\[
a_\mu^{\mathrm{seam}}
=
\frac54\cdot\frac{2304}{16777216\,\pi^8}\cdot\frac{1}{2\pi}
=
\frac{2880}{16777216\,\pi^9\cdot 2}
=
\frac{45}{524288\,\pi^9}.
\]
\end{proof}

\begin{remark}[Trace-theoretic reading]
The integer $2880=4!\,\tr_{S^+}(X^2)=24\times120$ with $X=6Y$ and
$\tr_{S^+}(X^2)=120=5!$, established in \cite[Corollary~5.9]{TFPT31},
encodes the quadratic hypercharge trace over the full positive half-spinor
$S^+=\Lambda^{\mathrm{even}}E$. The seam vertex correction therefore
carries the imprint of all $16$ chiral states in one family through the
factorization
\[
a_\mu^{\mathrm{seam}}
=
\frac{4!\,\tr_{S^+}(X^2)\,\cthree^8}{2\pi}.
\]
\end{remark}


\section{Comparison with experiment}\label{sec:comparison}

\subsection{Numerical evaluation}

\begin{theorem}[Numerical prediction]\label{thm:numerical-prediction}
The TFPT seam vertex correction evaluates to
\begin{equation}\label{eq:numerical-prediction}
a_\mu^{\mathrm{seam}}
=
\frac{\delta_2}{2\pi}
=
\frac{\num{1.80914597e-8}}{2\pi}
\approx
\num{2.879e-9}.
\end{equation}
\end{theorem}

\begin{proof}
From \Cref{lem:numerics}, $\delta_2\approx\num{1.80914597e-8}$.
Dividing by $2\pi\approx\num{6.28318531}$ yields
\[
a_\mu^{\mathrm{seam}}
\approx
\frac{\num{1.80914597e-8}}{\num{6.28318531}}
\approx
\num{2.8793e-9}.
\]
\end{proof}

\subsection{Comparison table}

\begin{table}[ht]
\centering
\begin{tabularx}{\linewidth}{@{}>{\raggedright\arraybackslash}p{0.32\linewidth}
    >{\centering\arraybackslash}p{0.30\linewidth}
    >{\centering\arraybackslash}X@{}}
\toprule
Source & Value $(\times10^{-9})$ & Status \\
\midrule
Experimental discrepancy $\Delta a_\mu$ & $2.49\pm0.48$ & Fermilab + BNL world average \\
TFPT seam vertex $\delta_2/(2\pi)$ & $2.879$ & exact from carrier algebra \\
\midrule
Deviation from central value & $+0.389$ & $0.81\sigma$ \\
\bottomrule
\end{tabularx}
\caption{The TFPT prediction versus the measured muon $g-2$ discrepancy. The
predicted value lies within $0.81\sigma$ of the experimental central value and
well inside the $1\sigma$ band.}
\label{tab:comparison}
\end{table}

The TFPT prediction $a_\mu^{\mathrm{seam}}\approx\num{2.879e-9}$ overshoots
the central experimental value by
\[
\frac{a_\mu^{\mathrm{seam}}-\Delta a_\mu^{\mathrm{central}}}{\sigma_{\Delta a_\mu}}
=
\frac{(\num{2.879}-\num{2.49})\times10^{-9}}{\num{0.48e-9}}
\approx
0.81,
\]
placing it well inside the $1\sigma$ experimental band.

\subsection{Parameter-free comparison}

The comparison has zero adjustable parameters:

\begin{enumerate}[label=(\roman*), leftmargin=1.6em, itemsep=0.25em]
\item The carrier split $E_3\oplus E_2$ is the unique solution of
      $\dim\Lambda^{\mathrm{even}}E=16$ and $\tr_EY^2=5/6$.
\item The invariants $B=3/2$, $\gammaY=5/6$, and $B\gammaY=5/4$ follow
      algebraically.
\item The seam constant $\cthree=1/(8\pi)$ follows from the unit spectral
      flow normalization.
\item The occupancy $\Oadm=48$ is the topological family closure output.
\item The defect coefficients $\deltatop$ and $\delta_2$ are derived from
      (ii)--(iv).
\item The projection factor $1/(2\pi)$ is the inverse seam-loop phase.
\end{enumerate}

No parameter is fitted to the muon data.

\section{Why \texorpdfstring{$\delta_2$}{delta2} and not
\texorpdfstring{$\deltatop$}{delta-top}}\label{sec:why-delta2}

The first-order defect $\deltatop$ enters the exact seam generating function
$\phiseam(\alpha)$, which determines the fine-structure constant $\alpha$
through the carrier-form closure equation
\[
\alpha^3-2\cthree^3\alpha^2-8\cthree^6\bone\ln\!\bigl((\phiseam(\alpha))^{-1}\bigr)=0.
\]
By construction, $\deltatop$ is \emph{consumed} by the electromagnetic closure:
it is part of the input that fixes $\alpha$.

The second-order defect $\delta_2=\frac54\deltatop^2$ is the leading term in the
seam opening beyond what the electromagnetic closure absorbs. In the Taylor
expansion
\[
\phiseam(\alpha)
=
\phibase+\deltatop e^{-2\alpha}+\delta_2\,e^{-4\alpha}+O(e^{-6\alpha}),
\]
the $\deltatop$ term feeds the $\alpha$-equation, while $\delta_2$ constitutes
the first residual correction that manifests as a separate observable effect.

In standard QFT language, $\deltatop$ corresponds to the charge
renormalization, while $\delta_2/(2\pi)$ corresponds to the anomalous magnetic
moment---the leading vertex correction beyond the coupling itself.

\begin{proposition}[Order assignment]
\label{prop:order-assignment}
The electromagnetic closure absorbs the first-order defect $\deltatop$ into
the determination of $\alpha$. The anomalous magnetic moment is the
leading residual topological observable and is therefore governed by the
second-order defect $\delta_2$.
\end{proposition}


\section{Discussion}\label{sec:discussion}

\subsection{No new particles}

The TFPT derivation does not invoke supersymmetric partners (sleptons,
charginos, neutralinos), dark photons, leptoquarks, or extended Higgs
sectors. The anomaly is not a signal of new particles propagating in
loops; it is the topological signature of the seam---the codimension-one
locus where the involution-doubled manifold $\widetilde M$ folds onto
the physical quotient $M$.

\subsection{Universality}

The vertex projection $\delta_2/(2\pi)$ is lepton-universal at the topological
level: the seam defect does not depend on which lepton traverses the vertex.
This is consistent with the fact that the electron anomalous magnetic
moment $a_e$ is dominated by the QED Schwinger series and is measured to a
precision where the $\sim\!\num{2.9e-9}$ topological correction would appear
at the level of the $\alpha^5$ QED coefficient, which is currently at the
frontier of both theory and experiment.

\subsection{Lattice HVP and the residual tension}

Recent lattice QCD evaluations of the hadronic vacuum polarization
(HVP)~\cite{BMW2021,RBC2024} suggest a smaller discrepancy than the
dispersive approach. If the lattice central value is adopted, the
experimental anomaly shrinks to $\Delta a_\mu\approx(1.5\pm0.9)\times10^{-9}$
(current estimates). The TFPT prediction $\num{2.879e-9}$ would then sit
at $\sim\!1.5\sigma$, still within the allowed band but less tightly centered.
The resolution of the HVP tension between dispersive and lattice methods is
an experimental question; the TFPT prediction is fixed and testable against
either outcome.

\subsection{Falsifiability}

The prediction $a_\mu^{\mathrm{seam}}=45/(524288\,\pi^9)$ is exact and
non-negotiable within TFPT. If future measurements of $\Delta a_\mu$ converge
to a value outside the range $\num{2.879e-9}\pm\num{0.5e-9}$, the seam
vertex mechanism in its present form would be excluded.

\section{Conclusion}\label{sec:conclusion}

The muon $g-2$ anomaly admits a parameter-free explanation within the
TFPT carrier algebra. The second-order topological defect
\[
\delta_2
=
\frac54\,\deltatop^2
=
\frac{45}{262144\,\pi^8}
\approx\num{1.809e-8},
\]
projected through the seam-loop holonomy phase $\Delta_\Seam=2\pi$, yields
\[
\boxed{
a_\mu^{\mathrm{seam}}
=
\frac{\delta_2}{2\pi}
=
\frac{45}{524288\,\pi^9}
\approx\num{2.879e-9}.
}
\]
This lies within $0.81\sigma$ of the Fermilab/BNL world average discrepancy
$\Delta a_\mu=(2.49\pm0.48)\times10^{-9}$ and requires no supersymmetry, no
dark photons, and no free parameters.

The muon anomaly, in this reading, is not a messenger of new particles. It
is the topological signature of the seam---the intrinsic codimension-one
defect of the carrier-algebraic spacetime geometry.

\appendix

\section{Full derivation chain}\label{app:chain}

For self-containedness, we collect the complete algebraic chain from the
carrier split to the muon anomaly in one sequence.

\medskip

\textbf{Step 1: Carrier.}
$E=E_3\oplus E_2$ with $\dim\Lambda^{\mathrm{even}}E=16$,\quad
$\gammaY=\tr_EY^2=5/6$,\quad $B=3/2$.

\textbf{Step 2: Compression quotient.}
$B\gammaY=5/4$.

\textbf{Step 3: Seam constant.}
$\SF(U_\Seam)=1$, $\Delta_\Seam=2\pi$,
$\cthree=\Delta_\Seam/(16\pi^2)=1/(8\pi)$.

\textbf{Step 4: Occupancy.}
$N_{\mathrm{fam}}=3$, $\dim S^+=16$,
$\Oadm=N_{\mathrm{fam}}\dim S^+=48$.

\textbf{Step 5: First-order defect.}
$\deltatop=\Oadm\,\cthree^4=48/(8\pi)^4=3/(256\pi^4)\approx\num{1.203e-4}$.

\textbf{Step 6: Second-order defect.}
$\delta_2=B\gammaY\,\deltatop^2=\frac54\,\deltatop^2
=45/(262144\,\pi^8)\approx\num{1.809e-8}$.

\textbf{Step 7: Vertex projection.}
$a_\mu^{\mathrm{seam}}=\delta_2/\Delta_\Seam=\delta_2/(2\pi)
=45/(524288\,\pi^9)\approx\num{2.879e-9}$.

\textbf{Step 8: Comparison.}
$\Delta a_\mu^{\exp}=(2.49\pm0.48)\times10^{-9}$.
Deviation: $0.81\sigma$.


\section{Numerical verification}\label{app:numerics}

Independent numerical check using extended precision:

\begin{align}
\pi &= \num{3.14159265358979324}\notag\\
\pi^4 &= \num{97.4090910340024}\notag\\
256\pi^4 &= \num{24936.7273}\notag\\
\deltatop &= 3/\num{24936.7273} = \num{1.20304484e-4}\notag\\
\deltatop^2 &= \num{1.44731677e-8}\notag\\
\delta_2 &= \tfrac54\times\num{1.44731677e-8} = \num{1.80914597e-8}\notag\\
2\pi &= \num{6.28318531}\notag\\
a_\mu^{\mathrm{seam}} &= \num{1.80914597e-8}/\num{6.28318531}
= \num{2.8793e-9}\label{eq:check}
\end{align}


\section{Imported TFPT kernel block}\label{app:import}

The minimal import block for this application note is
\[
\bigl(
E_3\oplus E_2,\,
Y,\,
\inv,\,
\gammaY,\,
B,\,
\cthree,\,
\deltatop,\,
\delta_2,\,
\Oadm
\bigr).
\]
\begin{sloppypar}
This block is intentionally small enough to be verified
independently against the parent manuscript~\cite{TFPT31}
and the clean kernel note~\cite{TFPTkernel}.
\end{sloppypar}

\begin{sloppypar}
\begin{thebibliography}{99}

\bibitem{TFPT31}
S.~Hamann and A.~Rizzo,
\emph{TFPT: Carrier Algebra, Closure Theorems, and Intrinsic Observation},
Version~4.1 (2025).

\bibitem{TFPTkernel}
S.~Hamann and A.~Rizzo,
\emph{TFPT Minimal Kernel: Clean core note for downstream
application papers} (2025).

\bibitem{Fermilab2021}
B.~Abi \emph{et al.} (Muon $g-2$ Collaboration),
\emph{Measurement of the Positive Muon Anomalous Magnetic Moment to 0.46 ppm},
Phys.\ Rev.\ Lett.\ \textbf{126}, 141801 (2021).

\bibitem{Fermilab2023}
D.~P.~Aguillard \emph{et al.} (Muon $g-2$ Collaboration),
\emph{Measurement of the Positive Muon Anomalous Magnetic Moment to 0.20 ppm},
Phys.\ Rev.\ Lett.\ \textbf{131}, 161802 (2023).

\bibitem{BNL2006}
G.~W.~Bennett \emph{et al.} (Muon $g-2$ Collaboration),
\emph{Final Report of the Muon E821 Anomalous Magnetic Moment Measurement at BNL},
Phys.\ Rev.\ D~\textbf{73}, 072003 (2006).

\bibitem{WP2020}
T.~Aoyama \emph{et al.},
\emph{The anomalous magnetic moment of the muon in the Standard Model},
Phys.\ Rep.\ \textbf{887}, 1 (2020).

\bibitem{BMW2021}
Sz.~Borsanyi \emph{et al.} (BMW Collaboration),
\emph{Leading hadronic contribution to the muon magnetic moment from lattice QCD},
Nature \textbf{593}, 51 (2021).

\bibitem{RBC2024}
T.~Blum \emph{et al.} (RBC/UKQCD Collaboration),
\emph{Update of the hadronic vacuum polarization contribution to $a_\mu$},
Phys.\ Rev.\ Lett.\ \textbf{132}, 051903 (2024).

\end{thebibliography}
\end{sloppypar}

\end{document}
